\documentclass[10pt,a4paper]{amsart}
\usepackage{amsmath,amssymb,amsthm,graphicx}
\usepackage{verbatim}
\newenvironment{reference}{\comment}{\endcomment}
%\newenvironment{reference}{}{}
\newenvironment{slogan}{\comment}{\endcomment}
\newenvironment{history}{\comment}{\endcomment}
% For commutative diagrams we use Xy-pic
\usepackage[all]{xy}

% We use 2cell for 2-commutative diagrams.
\xyoption{2cell}
% Theorem environments.
%
\theoremstyle{plain}
\newtheorem{theorem}[subsection]{Theorem}
\newtheorem{proposition}[subsection]{Proposition}
\newtheorem{lemma}[subsection]{Lemma}

\theoremstyle{definition}
\newtheorem{definition}[subsection]{Definition}
\newtheorem{example}[subsection]{Example}
\newtheorem{exercise}[subsection]{Exercise}
\newtheorem{situation}[subsection]{Situation}

\theoremstyle{remark}
\newtheorem{remark}[subsection]{Remark}
\newtheorem{remarks}[subsection]{Remarks}

\numberwithin{equation}{subsection}

% Macros
%
\def\lim{\mathop{\mathrm{lim}}\nolimits}
\def\colim{\mathop{\mathrm{colim}}\nolimits}
\def\Spec{\mathop{\mathrm{Spec}}}
\def\Hom{\mathop{\mathrm{Hom}}\nolimits}
\def\Ext{\mathop{\mathrm{Ext}}\nolimits}
\def\SheafHom{\mathop{\mathcal{H}\!\mathit{om}}\nolimits}
\def\SheafExt{\mathop{\mathcal{E}\!\mathit{xt}}\nolimits}
\def\Sch{\mathit{Sch}}
\def\Mor{\mathop{\mathrm{Mor}}\nolimits}
\def\Ob{\mathop{\mathrm{Ob}}\nolimits}
\def\Sh{\mathop{\mathit{Sh}}\nolimits}
\def\NL{\mathop{N\!L}\nolimits}
\def\CH{\mathop{\mathrm{CH}}\nolimits}
\def\proetale{{pro\text{-}\acute{e}tale}}
\def\etale{{\acute{e}tale}}
\def\QCoh{\mathit{QCoh}}
\def\Ker{\mathop{\mathrm{Ker}}}
\def\Im{\mathop{\mathrm{Im}}}
\def\Coker{\mathop{\mathrm{Coker}}}
\def\Coim{\mathop{\mathrm{Coim}}}

% Boxtimes
%
\DeclareMathSymbol{\boxtimes}{\mathbin}{AMSa}{"02}

%
% Macros for moduli stacks/spaces
%
\def\QCohstack{\mathcal{QC}\!\mathit{oh}}
\def\Cohstack{\mathcal{C}\!\mathit{oh}}
\def\Spacesstack{\mathcal{S}\!\mathit{paces}}
\def\Quotfunctor{\mathrm{Quot}}
\def\Hilbfunctor{\mathrm{Hilb}}
\def\Curvesstack{\mathcal{C}\!\mathit{urves}}
\def\Polarizedstack{\mathcal{P}\!\mathit{olarized}}
\def\Complexesstack{\mathcal{C}\!\mathit{omplexes}}
% \Pic is the operator that assigns to X its picard group, usage \Pic(X)
% \Picardstack_{X/B} denotes the Picard stack of X over B
% \Picardfunctor_{X/B} denotes the Picard functor of X over B
\def\Pic{\mathop{\mathrm{Pic}}\nolimits}
\def\Picardstack{\mathcal{P}\!\mathit{ic}}
\def\Picardfunctor{\mathrm{Pic}}
\def\Deformationcategory{\mathcal{D}\!\mathit{ef}}
\usepackage{cleveref}
\crefname{originalreference}{original reference}{original references}
\Crefname{originalreference}{Original reference}{Original references}
\makeatletter
\newlabel{algebra-lemma-apply-grothendieck-module}{{algebra-lemma-apply-grothendieck-module}{source}{}{sourcepointer-1}{}}
\newlabel{algebra-lemma-apply-grothendieck-module@cref}{{[originalreference][0][]algebra-lemma-apply-grothendieck-module}{[1][0][]source}{}{}{}}
\makeatother
\crefname{originalreference}{original reference}{original references}
\Crefname{originalreference}{Original reference}{Original references}
\makeatletter
\newlabel{dualizing-lemma-dualizing-gorenstein-formal-fibres}{{dualizing-lemma-dualizing-gorenstein-formal-fibres}{source}{}{sourcepointer-2}{}}
\newlabel{dualizing-lemma-dualizing-gorenstein-formal-fibres@cref}{{[originalreference][0][]dualizing-lemma-dualizing-gorenstein-formal-fibres}{[1][0][]source}{}{}{}}
\makeatother
\crefname{originalreference}{original reference}{original references}
\Crefname{originalreference}{Original reference}{Original references}
\makeatletter
\newlabel{more-algebra-section-properties-formal-fibres}{{more-algebra-section-properties-formal-fibres}{source}{}{sourcepointer-3}{}}
\newlabel{more-algebra-section-properties-formal-fibres@cref}{{[originalreference][0][]more-algebra-section-properties-formal-fibres}{[1][0][]source}{}{}{}}
\makeatother
\crefname{originalreference}{original reference}{original references}
\Crefname{originalreference}{Original reference}{Original references}
\makeatletter
\newlabel{more-algebra-lemma-completion-dimension}{{more-algebra-lemma-completion-dimension}{source}{}{sourcepointer-4}{}}
\newlabel{more-algebra-lemma-completion-dimension@cref}{{[originalreference][0][]more-algebra-lemma-completion-dimension}{[1][0][]source}{}{}{}}
\makeatother
\title{Comparison of depth and of the depth-dimension sum between a local ring and its ideal-adic completion}
\author{The Stacks Project Authors}
\date{stacks-project:local-cohomology.tex:local-snapshot-20261012; GNU FDL 1.2 or later}
\begin{document}
\maketitle
\noindent\textit{Editorial note.} All delivered mathematical text is the author's own, copied from the pinned Stacks Project chapter source sources/stacks-local-cohomology/pinned-original-source.tex; no statement or proof step was rewritten, reordered or removed.\par\medskip
\noindent\textit{Editorial note.} Preamble provenance: the Stacks Project keeps one single shared preamble for every chapter; this card takes it directly from the staged sources/stacks-local-cohomology/preamble.tex as three verbatim spans (lines 8-12, 14-18 and 154-217), so no derived file was written and nothing was added or redefined. The wrapper class and the base packages are supplied by the assembler.\par\medskip
\noindent\textit{Editorial note.} Bibliography provenance: neither the counted statement nor its proof contains a citation (this lane rescans the whole span before writing the card), so the card carries no bibliography and no entry was reconstructed.\par\medskip
\noindent\textit{Editorial note.} Selection provenance: this card was opened under the file-based control contract (CONTROL.md), with the mechanical pre-card gate run on its statement and proof interval before the card was written.\par\medskip
\noindent\textit{Allowed standard prerequisites (editorial).}\begin{itemize}
\item The behaviour of depth under a flat local homomorphism, together with the formal fibre formula. Invoked inside the counted proof at source lines 2585-2593 and attributed to a lemma of another chapter (a cross-chapter pointer). Splits the depth of the completion into the depth of the original module and the depth of the formal fibre.
\item The Cohen-Macaulay property of the formal fibres of a Noetherian ring with a dualizing complex, and the equality of dimensions of a local ring and its completion. Both are invoked inside the counted proof (source lines 2594-2601 and 2617-2620) as cross-chapter pointers, the first with a reference to a further chapter of more algebra. Replaces the depth of the formal fibre by its dimension and identifies the two dimension terms appearing in the displayed chain of equalities.
\end{itemize}\medskip
\ifdefined\hypertarget\hypertarget{sourcepointer-1}{}\fi
\noindent\textit{Source reference.} Original reference algebra-lemma-apply-grothendieck-module: sources/stacks-local-cohomology/pinned-original-source.tex, source lines 2593--2593. Navigation pointer only: the counted text refers to the Stacks item algebra-lemma-apply-grothendieck-module, which belongs to another chapter of the project and is not part of this card. The pointer carries the target's own label as its number because this chapter's staged snapshot does not contain the other chapter's tag table; its content is not reproduced here.\par\medskip
\ifdefined\hypertarget\hypertarget{sourcepointer-2}{}\fi
\noindent\textit{Source reference.} Original reference dualizing-lemma-dualizing-gorenstein-formal-fibres: sources/stacks-local-cohomology/pinned-original-source.tex, source lines 2596--2596. Navigation pointer only: the counted text refers to the Stacks item dualizing-lemma-dualizing-gorenstein-formal-fibres, which belongs to another chapter of the project and is not part of this card. The pointer carries the target's own label as its number because this chapter's staged snapshot does not contain the other chapter's tag table; its content is not reproduced here.\par\medskip
\ifdefined\hypertarget\hypertarget{sourcepointer-3}{}\fi
\noindent\textit{Source reference.} Original reference more-algebra-section-properties-formal-fibres: sources/stacks-local-cohomology/pinned-original-source.tex, source lines 2598--2598. Navigation pointer only: the counted text refers to the Stacks item more-algebra-section-properties-formal-fibres, which belongs to another chapter of the project and is not part of this card. The pointer carries the target's own label as its number because this chapter's staged snapshot does not contain the other chapter's tag table; its content is not reproduced here.\par\medskip
\ifdefined\hypertarget\hypertarget{sourcepointer-4}{}\fi
\noindent\textit{Source reference.} Original reference more-algebra-lemma-completion-dimension: sources/stacks-local-cohomology/pinned-original-source.tex, source lines 834--834. Navigation pointer only: the counted text refers to the Stacks item more-algebra-lemma-completion-dimension, which belongs to another chapter of the project and is not part of this card. The pointer carries the target's own label as its number because this chapter's staged snapshot does not contain the other chapter's tag table; its content is not reproduced here.\par\medskip
\begin{lemma}
\label{lemma-change-completion}
Let $A$ be a Noetherian ring which has a dualizing complex.
Let $I \subset A$ be an ideal.
Let $M$ be a finite $A$-module. Let $A', M'$ be the $I$-adic
completions of $A, M$. Let $\mathfrak p' \subset \mathfrak q'$
be prime ideals of $A'$ with $\mathfrak q' \in V(IA')$
lying over $\mathfrak p \subset \mathfrak q$ in $A$. Then
$$
\text{depth}_{A_{\mathfrak p'}}(M'_{\mathfrak p'})
\geq
\text{depth}_{A_\mathfrak p}(M_\mathfrak p)
$$
and
$$
\text{depth}_{A_{\mathfrak p'}}(M'_{\mathfrak p'}) +
\dim((A'/\mathfrak p')_{\mathfrak q'}) =
\text{depth}_{A_\mathfrak p}(M_\mathfrak p) +
\dim((A/\mathfrak p)_\mathfrak q)
$$
\end{lemma}

\begin{proof}
We have
$$
\text{depth}(M'_{\mathfrak p'}) =
\text{depth}(M_\mathfrak p) +
\text{depth}(A'_{\mathfrak p'}/\mathfrak p A'_{\mathfrak p'})
\geq \text{depth}(M_\mathfrak p)
$$
by flatness of $A \to A'$, see
Algebra, Lemma \ref{algebra-lemma-apply-grothendieck-module}.
Since the fibres of $A \to A'$ are Cohen-Macaulay
(Dualizing Complexes, Lemma
\ref{dualizing-lemma-dualizing-gorenstein-formal-fibres} and
More on Algebra, Section
\ref{more-algebra-section-properties-formal-fibres})
we see that
$\text{depth}(A'_{\mathfrak p'}/\mathfrak p A'_{\mathfrak p'}) =
\dim(A'_{\mathfrak p'}/\mathfrak p A'_{\mathfrak p'})$.
Thus we obtain
\begin{align*}
\text{depth}(M'_{\mathfrak p'}) +
\dim((A'/\mathfrak p')_{\mathfrak q'})
& =
\text{depth}(M_\mathfrak p) +
\dim(A'_{\mathfrak p'}/\mathfrak p A'_{\mathfrak p'}) +
\dim((A'/\mathfrak p')_{\mathfrak q'}) \\
& =
\text{depth}(M_\mathfrak p) +
\dim((A'/\mathfrak p A')_{\mathfrak q'}) \\
& =
\text{depth}(M_\mathfrak p) +
\dim((A/\mathfrak p)_\mathfrak q)
\end{align*}
Second equality because $A'$ is catenary and third equality by
More on Algebra, Lemma \ref{more-algebra-lemma-completion-dimension}
as $(A/\mathfrak p)_\mathfrak q$ and $(A'/\mathfrak p A')_{\mathfrak q'}$
have the same $I$-adic completions.
\end{proof}

\end{document}
